\documentclass{article}

\usepackage[utf8]{inputenc}
\usepackage[T1]{fontenc}
\usepackage{hyperref}
\usepackage{url}
\usepackage{booktabs}
\usepackage{amsmath}
\usepackage{microtype}
\usepackage{graphicx}

% Numbers are generated from a reproducible experiment, not typed in by hand:
%   make -C papers results        (writes papers/generated/multi_agent_*.tex)
% Each \ma... macro expands to a value produced by harness.run_experiment on the
% deterministic `sim` provider.
% AUTO-GENERATED by papers/generate_results.py -- do not edit by hand.
% Regenerate with: make -C papers results
\newcommand{\maSingleAccuracy}{1}
\newcommand{\maSingleTokensIn}{5.8}
\newcommand{\maSingleTokensOut}{4}
\newcommand{\maSelfConsistencyAccuracy}{1}
\newcommand{\maSelfConsistencyTokensIn}{28.8}
\newcommand{\maSelfConsistencyTokensOut}{20}
\newcommand{\maConsensusAccuracy}{1}
\newcommand{\maConsensusTokensIn}{179}
\newcommand{\maConsensusTokensOut}{20}
\newcommand{\maDebateAccuracy}{1}
\newcommand{\maDebateTokensIn}{81}
\newcommand{\maDebateTokensOut}{12}
\newcommand{\maManagerWorkerAccuracy}{1}
\newcommand{\maManagerWorkerTokensIn}{62}
\newcommand{\maManagerWorkerTokensOut}{12}


\title{The Price of Coordination: A Reproducible Benchmark for\\
Multi-Agent Compute Overhead}

\author{
  Hidden Layer Lab \\
  \texttt{contact@hiddenlayer.ai}
}

\date{}

\begin{document}

\maketitle

\begin{abstract}
Multi-agent strategies---debate, self-consistency, consensus, manager--worker---can
improve answer quality \cite{du2023debate,wang2023selfconsistency}, but they also do more
work: many model calls instead of one. The quality side of this trade-off depends on the
model under test; the \emph{cost} side does not, and is rarely reported in a reproducible
way. We present a small, provider-agnostic benchmark that isolates coordination overhead
by running every strategy through a deterministic simulated provider that fixes answer
quality. With accuracy held constant at \maSingleAccuracy{} across all five strategies,
input-token cost ranges from \maSingleTokensIn{} tokens for a single call to
\maConsensusTokensIn{} for consensus (Table~\ref{tab:ma-cost})---a $\sim$30$\times$ spread
attributable purely to coordination. Every number here is produced by the experiment
runner and injected via generated macros, demonstrating a reproducible
hypothesis$\rightarrow$test$\rightarrow$publish path that generalizes across projects.
\end{abstract}

\section{Introduction}
The case for multi-agent inference is a quality argument: debate and self-consistency
trade extra computation for better answers \cite{du2023debate,wang2023selfconsistency}.
Whether that trade pays off depends on the deployment---model, task, budget---so the
\emph{overhead} half of the equation deserves first-class, reproducible measurement.
Yet papers typically report quality deltas while leaving compute cost as an aside.

We separate the two axes. The quality axis is model-dependent and is studied with a live
provider (companion work). The \emph{cost axis}---how many tokens/calls a strategy spends
relative to a single call---is a property of the coordination structure itself, and can
be measured exactly and reproducibly by holding answer quality fixed. We do so with a
deterministic simulated provider, yielding a benchmark that runs offline, in CI, on any
machine, with identical results.

\section{Method}
We compare five strategies from the lab's multi-agent harness: \texttt{single}
(baseline), \texttt{self\_consistency}, \texttt{consensus}, \texttt{debate}, and
\texttt{manager\_worker}. Each is run over the same four short factual questions through
the harness's deterministic \texttt{sim} provider, which returns a fixed canned answer
for every call. Because the answer is fixed, every strategy scores identically on
keyword-match accuracy; the only thing that varies between strategies is how many model
calls they make and how large the prompts are. We report mean input and output tokens per
task as the overhead measure. The entire run is driven by a single experiment spec and
executed by \texttt{harness.run\_experiment}; the table and the macros cited in this text
are emitted by \texttt{harness.analysis}.

\section{Results}
Table~\ref{tab:ma-cost} reports accuracy and per-task token cost for each strategy. All
values are generated from the run.

% AUTO-GENERATED by papers/generate_results.py -- do not edit by hand.
% Regenerate with: make -C papers results
\begin{table}[t]
\centering
\begin{tabular}{lrrr}
\toprule
arm & accuracy & tokens\_in & tokens\_out \\
\midrule
single & 1 & 5.8 & 4 \\
self\_consistency & 1 & 28.8 & 20 \\
consensus & 1 & 179 & 20 \\
debate & 1 & 81 & 12 \\
manager\_worker & 1 & 62 & 12 \\
\bottomrule
\end{tabular}
\caption{Five coordination strategies on four factual items through the deterministic \texttt{sim} provider. Answer quality is held fixed (accuracy ties at 1.0), so the token counts isolate the \emph{computational overhead} of coordination relative to a single call. Real-model quality differences are out of scope here (they require a live provider).}
\label{tab:ma-cost}
\end{table}


With quality fixed (accuracy $=\maSingleAccuracy{}$ for every strategy), the overhead is
stark. Relative to the single-call baseline (\maSingleTokensIn{} input tokens),
self-consistency uses \maSelfConsistencyTokensIn{}, manager--worker \maManagerWorkerTokensIn{},
debate \maDebateTokensIn{}, and consensus \maConsensusTokensIn{}---roughly a thirty-fold
range driven entirely by the coordination structure (number of agents, rounds, and
synthesis steps). Output-token costs scale similarly, from \maSingleTokensOut{} to
\maConsensusTokensOut{}. These are exactly the multipliers a practitioner pays per query
before any quality benefit is realized.

\section{Discussion}
Reading this benchmark alongside a real-model quality study makes the trade-off concrete:
a strategy that lifts accuracy by a few points while costing $20$--$30\times$ the tokens
may or may not be worth it depending on budget and stakes, and now the cost term is a
measured number rather than a guess. Because the harness is provider-agnostic, swapping
\texttt{sim} for a live provider produces the quality axis on the same five arms with the
same code, so cost and quality can be reported in one table.

\paragraph{Reproducibility.} This paper is an artifact of the pipeline it describes:
\texttt{make -C papers results} re-runs the experiment and rewrites the macros and table
that this document \texttt{\textbackslash input}s, so the numbers cannot drift from the code.

\section{Limitations}
The simulated provider fixes answer content, so this benchmark deliberately says nothing
about answer \emph{quality}; it measures coordination overhead only. Token counts under
\texttt{sim} reflect the strategies' prompt-construction structure rather than a specific
tokenizer, so absolute values should be read as structural multipliers, not vendor token
bills. Latency is not reported because it is dominated by (here negligible) model time and
is not reproducible across machines.

\section{Conclusion}
Coordination is not free, and its cost is reproducibly measurable independent of any
model. We provide a provider-agnostic benchmark that isolates that cost, find a
$\sim$30$\times$ token spread across common strategies at fixed quality, and demonstrate a
numbers-from-runs publishing path that generalizes across the lab's projects.

\bibliographystyle{plain}
\bibliography{references}

\end{document}
